\documentclass[11pt]{article}
\usepackage{amsmath,amssymb,amsfonts}
\usepackage[letterpaper,margin=1in]{geometry}
\usepackage{array}

\title{Coupled Twiss Calculation in \texttt{elegant}}
\author{}
\date{}

\renewcommand{\arraystretch}{1.35}

\begin{document}
\maketitle

\section{Overview}

The \verb|coupled_twiss_output| command computes the equilibrium second
moments of a beam in a lattice with arbitrary linear coupling between the
transverse planes (and, optionally, the longitudinal plane). It follows the
Ripken / Mais formalism: rather than separate $\beta_x(s)$ and $\beta_y(s)$,
each normal mode $k$ of the lattice is described by a complex eigenvector
$\mathbf{v}_k$ of the one-turn linear map, and its contribution to the beam
moments at any element is captured by a real symmetric matrix $A_k(s)$ of
rank at most~$2$.

This note explains what the SDDS output columns mean, in particular the
cross-block entries \verb|A_xy_k|, \verb|A_xpy_k|, \verb|A_xyp_k|,
\verb|A_xpyp_k|, and the diagonal $\gamma$ entries, and describes how
\verb|matched=1| (periodic) and \verb|matched=0| (user-specified initial
conditions) propagate them along the beamline.

\section{Phase space and notation}

Phase space coordinates follow the standard \texttt{elegant} ordering
\[
   \mathbf{x} \;=\; \bigl( x,\;x',\;y,\;y',\;s,\;\delta \bigr)^T ,
\]
with indices $0,\dots,5$. The 6-D symplectic form is
\[
   S \;=\;
   \begin{pmatrix} J & & \\ & J & \\ & & J \end{pmatrix} ,
   \qquad
   J \;=\; \begin{pmatrix} 0 & 1 \\ -1 & 0 \end{pmatrix} .
\]
Calculations may be carried out in either 4-D (transverse only,
\verb|calculate_3d_coupling=0| or no cavity present) or 6-D (with synchrotron
motion). Let $M$ denote the one-turn linear transfer matrix and $M(s)$ the
accumulated matrix from the start of the calculation (the \verb|RECIRC|
element or the start of the beamline) to position $s$.

\section{Eigenvector decomposition}

For a stable linear map the eigenvalues of $M$ come in complex conjugate
pairs $\lambda_k = e^{\pm 2\pi i\nu_k}$, $k = 1,2$ (and $k=3$ in 6-D). Let
$\mathbf{v}_k \in \mathbb{C}^N$ be the right eigenvector of the pair
representative,
\[
   M\,\mathbf{v}_k \;=\; e^{2\pi i \nu_k}\,\mathbf{v}_k ,
\]
and write $\mathbf{v}_k = \mathbf{e}_k + i\,\mathbf{f}_k$ with
$\mathbf{e}_k,\mathbf{f}_k \in \mathbb{R}^N$. Symplectic normalization is
chosen so that
\[
   \bigl|\,\mathbf{v}_k^\dagger S\,\mathbf{v}_k\,\bigr| \;=\; 1
   \qquad \text{for every $k$.}
\]
After diagonalization the conjugate pairs are sorted so that mode~1 has its
dominant projection in the $(x,x')$ subspace, mode~2 in $(y,y')$, and
mode~3 in $(s,\delta)$ -- the same convention used in the uncoupled limit.

\section{The mode matrices $A_k$}

For each mode define the real symmetric matrix
\[
   \boxed{\;
   A_k \;=\; \mathbf{e}_k\,\mathbf{e}_k^T \;+\; \mathbf{f}_k\,\mathbf{f}_k^T .
   \;}
\]
$A_k$ is positive semi-definite, of rank at most $2$, and propagates exactly
like the second-moment matrix:
\[
   \boxed{\; A_k(s) \;=\; M(s)\, A_k(0)\, M(s)^T . \;}
\]
This follows from the linearity of the transport: each component
$\mathbf{e}_k$, $\mathbf{f}_k$ transports as $M(s)\,\mathbf{e}_k(0)$, so the
outer-product structure of $A_k$ is preserved.

\paragraph{Connection to the sigma matrix.}
Given mode emittances $\varepsilon_k$ (for the transverse modes these come
from the namelist or the \verb|twiss_output| radiation integrals; for the
longitudinal mode $\varepsilon_3$ is computed from $\sigma_\delta$), the
beam covariance matrix is
\[
   \Sigma(s) \;=\; \sum_k \varepsilon_k\, A_k(s) ,
   \qquad
   \langle x_i\,x_j\rangle \;=\; \sum_k \varepsilon_k\, (A_k)_{ij} .
\]
So $A_k$ is, up to an emittance factor, the contribution of mode~$k$ to the
beam covariance.

\section{Per-plane Twiss-like parameters}

The diagonal $2\times 2$ blocks of $A_k$ in the transverse planes give the
generalized Twiss functions of mode~$k$:
\begin{center}
\begin{tabular}{l l l}
\hline
quantity & matrix entry & SDDS column \\
\hline
$\beta_x^{(k)}$         & $(A_k)_{xx}$    & \verb|betax|$k$  \\
$\gamma_x^{(k)}$        & $(A_k)_{x'x'}$  & \verb|gammax|$k$ \\
$-\alpha_x^{(k)}$       & $(A_k)_{xx'}$   & \verb|alphax|$k$ (stored with sign) \\
$\beta_y^{(k)}$         & $(A_k)_{yy}$    & \verb|betay|$k$  \\
$\gamma_y^{(k)}$        & $(A_k)_{y'y'}$  & \verb|gammay|$k$ \\
$-\alpha_y^{(k)}$       & $(A_k)_{yy'}$   & \verb|alphay|$k$ (stored with sign) \\
\hline
\end{tabular}
\end{center}
The naming convention follows the column header: e.g.\ \verb|betax2| is
the contribution of mode~2 (the $y$-like mode) to the $x$-plane beta function.
In an uncoupled lattice $\beta_x^{(2)} = \gamma_x^{(2)} = \alpha_x^{(2)} = 0$
and the diagonal block of $A_2$ in $(x,x')$ vanishes.

\paragraph{Twiss invariant in the coupled case.}
For an uncoupled single-plane Twiss eigenvector
$\mathbf{v} = (\sqrt{\beta},\,-(\alpha+i)/\sqrt{\beta})$ one verifies
$\beta\gamma - \alpha^2 = 1$. In a coupled lattice, however, mode~$k$'s
energy is shared between the two transverse planes, and the rank-$2$
constraint on $A_k$ applies to the full $4\times 4$ matrix rather than to
each plane separately. The plane sub-determinants then satisfy
\[
   \beta_x^{(k)}\,\gamma_x^{(k)} - \bigl(\alpha_x^{(k)}\bigr)^2 \;\le\; 1 ,
   \qquad
   \beta_y^{(k)}\,\gamma_y^{(k)} - \bigl(\alpha_y^{(k)}\bigr)^2 \;\le\; 1 ,
\]
with equality only when the mode is fully concentrated in that plane.
Consequently $\gamma_y^{(1)}$ is \emph{not} determined by $(\beta_y^{(1)},
\alpha_y^{(1)})$ once the mode also has support in the $x$-plane; it is an
independent quantity and is reported as a separate output column.

\section{Cross-block entries: \texttt{A\_xy\_k}, etc.}

Beyond the two diagonal $2\times 2$ blocks, each $A_k$ has a $2\times 2$
cross-block coupling $(x,x')$ to $(y,y')$. By symmetry it has four
independent real entries, reported as

\[
\begin{array}{rcl}
\verb|A_xy_|k   &=& (A_k)_{xy}   \;=\; \langle x\,y \rangle_k / \varepsilon_k , \\[2pt]
\verb|A_xpy_|k  &=& (A_k)_{x'y}  \;=\; \langle x'\,y \rangle_k / \varepsilon_k , \\[2pt]
\verb|A_xyp_|k  &=& (A_k)_{xy'}  \;=\; \langle x\,y' \rangle_k / \varepsilon_k , \\[2pt]
\verb|A_xpyp_|k &=& (A_k)_{x'y'} \;=\; \langle x'\,y' \rangle_k / \varepsilon_k ,
\end{array}
\]

where $\langle\cdot\rangle_k$ denotes the contribution of mode~$k$ alone to the
covariance. Together with $\beta,\gamma,\alpha$ for each plane these
fully specify the transverse $4\times 4$ $A_k$:
\[
A_k \;=\;
\left(
\begin{array}{cc|cc}
\beta_x^{(k)}      & -\alpha_x^{(k)}     & \verb|A_xy_|k     & \verb|A_xyp_|k \\
-\alpha_x^{(k)}    & \gamma_x^{(k)}      & \verb|A_xpy_|k    & \verb|A_xpyp_|k \\
\hline
\verb|A_xy_|k      & \verb|A_xpy_|k      & \beta_y^{(k)}     & -\alpha_y^{(k)} \\
\verb|A_xyp_|k     & \verb|A_xpyp_|k     & -\alpha_y^{(k)}   & \gamma_y^{(k)}
\end{array}
\right) .
\]
The cross-block vanishes in an uncoupled lattice. It grows with the
strength of coupling and, like the diagonals, is propagated by
$M(s)\,A_k(0)\,M(s)^T$.

\section{Beam sizes and tilt}

From the diagonal entries of $\Sigma$ the reported beam sizes are
\[
   \sigma_x = \sqrt{\Sigma_{xx}}, \quad
   \sigma_{x'} = \sqrt{\Sigma_{x'x'}}, \quad
   \sigma_y = \sqrt{\Sigma_{yy}}, \quad
   \sigma_{y'} = \sqrt{\Sigma_{y'y'}},
\]
and (in 6-D) the bunch length $\sigma_s = \sqrt{\Sigma_{ss}}$. The reported
$xy$ tilt angle is the principal-axis rotation of the $(x,y)$ ellipse,
\[
   \theta_{xy} \;=\;
   \tfrac{1}{2}\,\arctan\!\left(
     \frac{2\,\Sigma_{xy}}{\Sigma_{xx} - \Sigma_{yy}}
   \right) .
\]

\section{Dispersion}

\paragraph{Periodic case.}
In the 6-D matched case the dispersion is read off from the longitudinal-mode
matrix $A_3$. The eigenvector $\mathbf{v}_3$ has the closed-orbit-plus-
synchrotron structure
\[
   \mathbf{v}_3 \;=\;
   \bigl( \eta_x\,v_\delta,\; \eta_{x'} v_\delta,\;
          \eta_y\,v_\delta,\; \eta_{y'} v_\delta,\;
          v_s,\; v_\delta \bigr) ,
\]
where $(v_s, v_\delta)$ is the purely longitudinal piece. Then
\[
   (A_3)_{i\delta} \;=\; \eta_i\,|v_\delta|^2 \;=\; \eta_i\,(A_3)_{\delta\delta} ,
\]
and the signed dispersions are
\[
   \boxed{\;
   \eta_i \;=\; \frac{(A_3)_{i\delta}}{(A_3)_{\delta\delta}} ,
   \qquad i \in \{x,\,x',\,y,\,y'\} .
   \;}
\]
These four values are reported as \verb|etax|, \verb|etaxp|, \verb|etay|,
\verb|etayp|. They are signed and exact; an older formula
$\eta_x = \sqrt{(A_3)_{xx}\,(A_3)_{ss}}$ yields $|\eta_x|\sqrt{1+\alpha_s^2}$,
which differs from this by a small factor whenever $\alpha_s\ne 0$.

In a 4-D matched calculation no longitudinal mode is available and the
four eta columns are reported as zero.

\paragraph{Non-periodic case.}
With \verb|matched=0| the dispersion is propagated as a separate 6-vector
\[
   \mathbf{d}(s) \;=\; M(s)\,\mathbf{d}(0) ,
   \qquad
   \mathbf{d}(0) \;=\;
   \bigl(\eta_x^0,\; \eta_{x'}^0,\; \eta_y^0,\; \eta_{y'}^0,\; 0,\; 1\bigr)^T ,
\]
with the initial components supplied through the namelist parameters
\verb|eta_x|, \verb|etap_x|, \verb|eta_y|, \verb|etap_y|. The reported
$\eta$ columns are the first four components of $\mathbf{d}(s)$.

\section{The two operating modes}

\subsection*{\texttt{matched = 1} (default)}

\begin{enumerate}
\item \verb|accumulate_matrices| computes $M$ (one-turn) and side-effects
      $M(s)$ into \verb|eptr->accumMatrix| for every element.
\item LAPACK \verb|dgeev| diagonalizes $M$. Eigenvalues give the fractional
      tunes $\nu_x,\nu_y$ (reported as parameters \verb|nux|, \verb|nuy|);
      eigenvectors are symplectically normalized and sorted by dominant plane.
\item For each element, $A_k(s)$ is formed from the eigenvectors transported
      to $s$, the sigma matrix is assembled, and all output columns are
      written.
\end{enumerate}

\subsection*{\texttt{matched = 0}}

The user supplies the initial $A_k(0)$ matrices (and the initial dispersion
vector) via the namelist. The full set of inputs is
\[
\begin{aligned}
&\verb|beta_x1, beta_x2, beta_y1, beta_y2|, \\
&\verb|alpha_x1, alpha_x2, alpha_y1, alpha_y2|, \\
&\verb|gamma_x1, gamma_x2, gamma_y1, gamma_y2|, \\
&\verb|A_xy_1, A_xpy_1, A_xyp_1, A_xpyp_1|, \\
&\verb|A_xy_2, A_xpy_2, A_xyp_2, A_xpyp_2|, \\
&\verb|eta_x, etap_x, eta_y, etap_y| ,
\end{aligned}
\]
which exactly mirrors the SDDS output column names. The two $A_k(0)$
matrices are filled in directly from these values, the dispersion vector
$\mathbf{d}(0)$ is loaded, and both are propagated as in the matched case.

If \verb|gamma_x|$k$ or \verb|gamma_y|$k$ is left at its default value $-1$,
the code falls back to a rank-2 eigenvector heuristic,
\begin{align*}
   \gamma\big|_{\text{primary plane}}   &= (1+\alpha^2)/\beta , \\
   \gamma\big|_{\text{secondary plane}} &= \alpha^2/\beta ,
\end{align*}
where mode~1's primary plane is $x$ and mode~2's is $y$. The heuristic is
exact for an uncoupled initial condition; for a coupled one it is only an
approximation. Likewise, the cross-block entries default to zero
(block-diagonal $A_k$). For an exact round-trip of a periodic solution the
user should load \emph{all} the new columns from the reference \verb|.ctwi|
file rather than relying on the defaults.

\section{Round-tripping a periodic solution}

A common workflow is to compute a periodic optic, then study how it would
propagate under a different set of errors or initial conditions. To make
the non-periodic run reproduce the periodic one exactly, save the matched
\verb|.ctwi| file with \verb|matched=1| and then in the \verb|matched=0|
input deck load all four sets of initial conditions through \verb|rpn_load|:
\begin{verbatim}
&rpn_load filename = "ref.ctwi", tag = cref &end
&coupled_twiss_output
    matched   = 0,
    beta_x1   = "(cref.betax1)",   beta_x2 = "(cref.betax2)",
    beta_y1   = "(cref.betay1)",   beta_y2 = "(cref.betay2)",
    alpha_x1  = "(cref.alphax1)",  ...
    gamma_x1  = "(cref.gammax1)",  ...
    A_xy_1    = "(cref.A_xy_1)",   ...
    eta_x     = "(cref.etax)",     etap_x = "(cref.etaxp)",
    eta_y     = "(cref.etay)",     etap_y = "(cref.etayp)",
&end
\end{verbatim}
With all $24$ initial values supplied the propagated output agrees with the
matched output to within matrix-concatenation round-off: typically
$\sim 10^{-4}$ relative for the transverse betas and
$\sim 10^{-14}$ for the signed dispersions.

\section{Summary}

The output of \verb|coupled_twiss_output| reports, for each transverse mode
$k\in\{1,2\}$ at each element:

\begin{center}
\begin{tabular}{ll}
\hline
column & meaning \\
\hline
\verb|betax|$k$, \verb|betay|$k$    & diagonal: $(A_k)_{xx},(A_k)_{yy}$ \\
\verb|gammax|$k$, \verb|gammay|$k$  & diagonal: $(A_k)_{x'x'},(A_k)_{y'y'}$ \\
\verb|alphax|$k$, \verb|alphay|$k$  & sign-flipped diagonal: $-(A_k)_{xx'},-(A_k)_{yy'}$ \\
\verb|A_xy_|$k$, \verb|A_xpy_|$k$,
\verb|A_xyp_|$k$, \verb|A_xpyp_|$k$ & the four cross-block entries of $A_k$ \\
\verb|etax|, \verb|etaxp|, \verb|etay|, \verb|etayp| & signed dispersion components (6-D only) \\
\verb|Sx|, \verb|Sxp|, \verb|Sy|, \verb|Syp|, \verb|Ss|, \verb|xyTilt| & beam sizes and tilt from $\Sigma$ \\
\verb|S|$ij$ (if \verb|output_sigma_matrix=1|) & the full sigma matrix \\
\hline
\end{tabular}
\end{center}

These quantities, together with the propagation law
$A_k(s)=M(s)\,A_k(0)\,M(s)^T$ and the separate dispersion-vector transport,
are sufficient to reconstruct the full coupled second-moment matrix
$\Sigma(s)$ at every element.

\end{document}
